\chapter{Pathology}

\medterm{Acute Inflammation}-- A rapid, predominantly non-specific immune response to tissue injury or infection that typically lasts hours to days and involves vascular and cellular reactions. The cardinal signs are redness, heat, swelling, pain, and loss of function, resulting from increased blood flow, vascular permeability, and cellular infiltration. The process begins with vasodilation mediated by histamine, prostaglandins, and nitric oxide, followed by increased vascular permeability allowing plasma proteins and leukocytes to exit the circulation. Neutrophils are the first leukocytes recruited, followed by monocytes that transform into tissue macrophages, working together to eliminate the inciting stimulus and initiate tissue repair.

\medterm{Adrenal Cortical Carcinoma}-- A rare, aggressive malignancy arising from the adrenal cortex, with an incidence of approximately one to two per million population per year. The tumor may be functional, producing excessive hormones, or non-functional. Functional tumors may produce cortisol causing Cushing syndrome, aldosterone causing Conn syndrome, or androgens causing virilization in women or feminization in men. Non-functional tumors present with abdominal pain, early satiety, or are discovered incidentally on abdominal imaging. Complete surgical resection is the treatment of choice, but prognosis is poor, with five-year survival rates under 40 percent due to advanced stage at presentation and high recurrence rates.

\medterm{Adult Still's Disease}-- Adult Still's disease, also known as systemic-onset juvenile idiopathic arthritis when occurring in children, is a systemic inflammatory disorder characterized by high-spiking daily fevers, evanescent salmon-colored rash, arthritis, and multiorgan involvement. The pathogenesis involves excessive activation of the innate immune system with overproduction of IL-1, IL-6, IL-18, and TNF. The characteristic fever pattern is quotidian, meaning it occurs once daily typically in the evening, with rapid return to baseline between spikes. The rash is typically salmon-coloured, transient, and associated with fever spikes. Arthritis may be severe and destructive. Laboratory findings include markedly elevated ferritin levels, leukocytosis, and elevated inflammatory markers. Treatment involves NSAIDs, corticosteroids, and biologic agents targeting IL-1 or IL-6.

\medterm{Air Embolism}-- See forensic_medicine.

\medterm{Amyloidosis}-- A group of disorders characterized by the extracellular deposition of misfolded proteins that polymerize into insoluble fibrils with a distinctive beta-pleated sheet configuration. These fibrils bind Congo red dye and show characteristic apple-green birefringence under polarized light microscopy. AA amyloidosis results from chronic inflammatory conditions with serum amyloid A protein depositing predominantly in the kidneys, spleen, and liver. AL amyloidosis involves immunoglobulin light chain deposition from plasma cell dyscrasias, affecting the heart, kidneys, and peripheral nerves. Diagnosis requires tissue biopsy with Congo red staining showing apple-green birefringence under polarised light. Treatment targets the underlying plasma cell disorder.

\medterm{Anaphylaxis}-- See emergency_medicine.

\medterm{Anaplasia}-- A lack of differentiation, describing cells that have reverted to a more primitive, embryonic phenotype with loss of the specialized structural and functional characteristics of their tissue of origin. Anaplastic cells exhibit marked pleomorphism with variation in cell size and shape, hyperchromatic nuclei with coarse irregular chromatin, prominent and irregular nucleoli, and a high nuclear-to-cytoplasmic ratio exceeding 80 percent. Mitotic figures are abundant, frequently atypical, including tripolar and quadripolar forms. Anaplasia is a hallmark of malignancy and correlates with aggressive clinical behaviour, high histological grade, and poor prognosis. The most anaplastic tumors are classified as poorly differentiated or undifferentiated (grade IV).

\medterm{Angiogenesis}-- See oncology.

\medterm{APC Tumor Suppressor}-- The adenomatous polyposis coli protein is a large multidomain tumor suppressor that functions as a critical negative regulator of the canonical Wnt signaling pathway. In the absence of Wnt ligands, APC forms a destruction complex with axin, glycogen synthase kinase-3 beta, and casein kinase 1, which phosphorylates beta-catenin, targeting it for ubiquitination and proteasomal degradation. When APC is inactivated by mutation, the destruction complex cannot degrade beta-catenin, which accumulates in the nucleus and activates transcription of pro-proliferative genes including MYC and cyclin D1. Germline mutations in APC cause familial adenomatous polyposis, with hundreds of colorectal adenomas and near-100 percent risk of colorectal cancer by age 40.

\medterm{Apoptosis}-- See oncology.

\medterm{Apoptosis Extrinsic Pathway}-- The extrinsic pathway of apoptosis is triggered by extracellular ligands binding to death receptors belonging to the tumor necrosis factor receptor superfamily, including Fas (CD95), TNFR1, and TRAIL receptors DR4 and DR5. Upon ligand binding, these receptors trimerize and recruit intracellular adaptor proteins including FADD and TRADD through their death domains, forming the death-inducing signaling complex. FADD then recruits procaspase-8 through death effector domain interactions, leading to autocatalytic activation of caspase-8, which then cleaves and activates effector caspases-3 and -7. The extrinsic pathway can also engage the intrinsic pathway through BID cleavage, amplifying the apoptotic signal. This pathway is critical for immune surveillance and lymphocyte homeostasis.

\medterm{Apoptosis Intrinsic Pathway}-- The intrinsic, or mitochondrial, pathway of apoptosis is initiated by intracellular stress signals including DNA damage, growth factor withdrawal, endoplasmic reticulum stress, and hypoxia. These signals activate BH3-only proteins such as BIM, BAD, BID, and PUMA, which neutralize anti-apoptotic BCL-2 family members or directly activate the pro-apoptotic effectors BAX and BAK. Activated BAX and BAK oligomerize and form pores in the outer mitochondrial membrane, a process called mitochondrial outer membrane permeabilisation, releasing cytochrome c and other pro-apoptotic factors. Cytochrome c binds APAF-1 to form the apoptosome, activating caspase-9. This pathway is tightly regulated by the BCL-2 family and is disrupted in many cancers, contributing to chemotherapy resistance.

\medterm{Atherosclerotic Plaque}-- An atherosclerotic plaque is a lesion within the arterial intima composed of a necrotic lipid core covered by a fibrous cap, representing the pathological substrate of atherosclerotic cardiovascular disease. The necrotic core contains cholesterol crystals, calcium deposits, cellular debris, and foam cells derived from lipid-laden macrophages. The fibrous cap consists of smooth muscle cells, collagen, elastin, and proteoglycans that separate the thrombogenic core from the flowing blood. The plaque may be stable, with a thick fibrous cap and small lipid core, or vulnerable, with a thin cap, large lipid core, and dense inflammation. Rupture of a vulnerable plaque exposes thrombogenic core material, triggering acute thrombosis responsible for myocardial infarction and stroke.

\medterm{Autoimmune Disease}-- Autoimmune diseases result from loss of immunological self-tolerance, leading to immune-mediated attack on the body's own tissues. The mechanisms include molecular mimicry, where microbial antigens share epitopes with self-antigens triggering cross-reactive immune responses; bystander activation, where tissue damage releases sequestered self-antigens; epitope spreading, where immune responses extend to additional self-epitopes; and defective regulatory T cell function. Systemic lupus erythematosus, rheumatoid arthritis, type 1 diabetes, multiple sclerosis, and inflammatory bowel disease are common examples. Treatment typically involves immunosuppressive agents including corticosteroids, disease-modifying anti-rheumatic drugs, and biologic therapies targeting specific immune pathways.

\medterm{Bacterial Endocarditis Vegetations}-- Vegetations in bacterial endocarditis are complex masses composed of bacteria, fibrin, and platelets that form on heart valve surfaces or endocardium. The pathogenesis begins with bacteremia seeding a site of endothelial damage or a previously abnormal valve surface. The damaged endothelium exposes subendothelial collagen and von Willebrand factor, promoting platelet adhesion and fibrin deposition, creating a non-bacterial thrombotic endocarditis nidus. Bacteria, particularly Staphylococcus aureus, Streptococcus viridans, and Enterococcus species, adhere to the thrombotic nidus and proliferate, embedded in layers of fibrin and platelets. The vegetations can embolise to distant organs and destroy valve tissue, leading to acute or subacute endocarditis with fever, heart murmur, and peripheral stigmata.

\medterm{Behcet Disease}-- A chronic, relapsing systemic vasculitis characterized by recurrent oral ulcers, genital ulcers, and ocular inflammation, with a distribution along the ancient Silk Road from the Mediterranean to East Asia. The pathogenesis involves inappropriate neutrophil and T-cell mediated vascular inflammation, with genetic susceptibility associated with HLA-B51. The hallmark oral aphthous ulcers are painful, recurrent, and heal without scarring. Genital ulcers typically affect the scrotum and vulva, healing with scarring. Ocular involvement manifests as anterior or posterior uveitis and retinal vasculitis, potentially causing blindness. Neurological and large-vessel involvement carries the highest morbidity. Diagnosis is clinical. Treatment includes corticosteroids, immunosuppressants, and biologic agents.

\medterm{Benign Neoplasm}-- Benign neoplasms are proliferations of cells that remain confined to their site of origin, grow by expansion rather than invasion, and do not metastasize to distant sites. They typically grow slowly, pushing aside adjacent normal tissue and often becoming enclosed by a fibrous pseudocapsule formed from compressed connective tissue. Histologically, benign neoplasms show well-differentiated cells with regular nuclear morphology, low mitotic activity, no atypical mitotic figures, and growth pattern that preserves normal tissue architecture. They are usually encapsulated or well-circumscribed. Clinical effects result from local compression, hormone production, or obstruction of adjacent structures. Complete surgical excision is generally curative with an excellent prognosis.

\medterm{Bilirubin Metabolism}-- Bilirubin is the breakdown product of heme from senescent red blood cells, metabolized primarily by the liver through a series of enzymatic steps. Heme is converted to biliverdin by heme oxygenase, then to unconjugated bilirubin by biliverdin reductase. Unconjugated bilirubin is lipophilic and transported bound to albumin to the liver. In hepatocytes, unconjugated bilirubin is conjugated with glucuronic acid by UDP-glucuronosyltransferase, forming water-soluble conjugated bilirubin that is excreted into bile. In the intestine, conjugated bilirubin is deconjugated by bacterial enzymes and reduced to urobilinogen. Most urobilinogen is excreted in feces, but a portion undergoes enterohepatic circulation and is excreted in urine. Elevation of unconjugated or conjugated bilirubin causes clinical jaundice.

\medterm{Biopsy}-- See general.

\medterm{Carcinoid Tumor}-- Carcinoid tumors are well-differentiated neuroendocrine neoplasms that arise from enterochromaffin cells throughout the gastrointestinal tract and bronchial system. The most common primary sites are the appendix, small intestine, and rectum. The tumors produce biogenic amines and peptides, but clinically significant hormone secretion occurs mainly with hepatic metastases, which allow vasoactive substances to bypass hepatic first-pass metabolism. Carcinoid syndrome occurs in approximately 10 percent of patients, presenting with flushing, diarrhea, wheezing, and right-sided valvular heart disease from serotonin and other vasoactive substances. Diagnosis involves urinary 5-HIAA measurement. Surgical resection is curative for localised disease; somatostatin analogues control symptoms in metastatic disease.

\medterm{Carcinoma}-- See oncology.

\medterm{Cardiomyopathy}-- Diseases of the myocardium that result in cardiac dysfunction, excluding coronary artery disease, valvular disease, and congenital heart disease as primary causes. The three main types are dilated, hypertrophic, and restrictive cardiomyopathy. Dilated cardiomyopathy is characterized by ventricular dilation and impaired systolic function with reduced ejection fraction. It may be idiopathic, familial with mutations in genes encoding cytoskeletal and sarcomeric proteins, or secondary to causes such as alcoholic cardiomyopathy, chemotherapy-induced cardiotoxicity, or viral myocarditis. Hypertrophic cardiomyopathy is characterized by left ventricular hypertrophy without dilation, typically due to sarcomeric protein mutations. Restrictive cardiomyopathy involves impaired diastolic filling with preserved systolic function.

\medterm{Caseous Necrosis}-- A distinctive form of cell death that characteristically occurs in granulomatous infections, most notably tuberculosis. The term derives from the cheese-like, friable, whitish-yellow appearance of the necrotic tissue on gross examination. Histologically, caseous necrosis appears as amorphous, granular, eosinophilic debris without recognizable cellular outlines or tissue architecture, distinct from the structural preservation seen in coagulative necrosis. The process results from the combination of macrophage activation, T-cell mediated immunity, and bacterial products creating a hypoxic microenvironment. The necrotic material contains dead immune cells and bacterial debris. It may undergo liquefaction and cavitation, facilitating spread of infection to other sites.

\medterm{Cellular Adaptation}-- Cellular adaptations are reversible changes in cell morphology and function that occur in response to altered physiological demands or mild pathological conditions. Hypertrophy involves an increase in cell size through increased protein synthesis and organelle production, resulting in increased organ size, as seen in cardiac hypertrophy from chronic hypertension or skeletal muscle growth with resistance exercise. Hyperplasia involves an increase in cell number through mitotic division, as seen in the breast and endometrium in response to hormonal stimulation. Atrophy involves decreased cell size and organ mass due to reduced functional demand, inadequate blood supply, or ageing. Metaplasia is a reversible change where one differentiated cell type is replaced by another more resilient to environmental stress.

\medterm{Cerebral Edema}-- The accumulation of excess fluid within the brain parenchyma, leading to increased intracranial pressure and potentially fatal brain herniation. There are three main types with different pathophysiological mechanisms. Vasogenic edema results from disruption of the blood-brain barrier, allowing protein-rich fluid to enter the extracellular space; it is seen in brain tumors, abscesses, hypertensive encephalopathy, and posterior reversible encephalopathy syndrome. Cytotoxic edema results from energy failure and impaired ion pumps, causing intracellular swelling; it is seen in cerebral ischemia, hypoxia, and metabolic disorders. Interstitial edema occurs in obstructive hydrocephalus from transependymal CSF flow. Management focuses on reducing intracranial pressure through osmotic therapy, corticosteroids, or surgical decompression.

\medterm{Cerebral Infarction}-- Cerebral infarction, or ischemic stroke, results from occlusion of a cerebral artery causing irreversible neuronal damage. The most common cause is thromboembolism from atherosclerotic carotid or intracranial arteries, or cardioembolism from atrial fibrillation or cardiac thrombus. The brain is particularly vulnerable to ischemia due to its high metabolic rate and limited energy reserves, with neurons dying within minutes of complete ischemia. The ischemic core is surrounded by a penumbra of potentially salvageable tissue where blood flow is reduced but above the threshold for cell death. Clinical presentation depends on the affected vascular territory. Treatment includes thrombolysis within the therapeutic window, endovascular thrombectomy, and secondary prevention with antiplatelet agents and risk factor modification.

\medterm{Cholangiocarcinoma}-- A malignant tumor arising from the biliary epithelium, classified based on anatomical location as intrahepatic, perihilar, or distal extrahepatic. Risk factors include primary sclerosing cholangitis, particularly in association with inflammatory bowel disease, parasitic liver fluke infections, hepatolithiasis, and congenital biliary anomalies including choledochal cysts. Intrahepatic cholangiocarcinoma presents as a liver mass with abdominal pain and weight loss. Perihilar (Klatskin) tumors are the most common type and present with obstructive jaundice. Distal extrahepatic tumors have the best resectability rate. Overall prognosis is poor, with five-year survival under 20 percent. Surgical resection with negative margins offers the only chance of cure; liver transplantation may be considered in selected cases.

\medterm{Cholecystitis}-- Inflammation of the gallbladder, most commonly caused by obstruction of the cystic duct by a gallstone, which occurs in approximately 95 percent of cases. The obstruction leads to bile stasis, increased intraluminal pressure, and ischemia of the gallbladder wall, which may progress to bacterial infection. Acute cholecystitis presents with right upper quadrant pain, fever, nausea, and vomiting, with physical findings including Murphy's sign, arrest of inspiration during palpation of the right upper quadrant. Ultrasound findings include gallbladder wall thickening, distension, pericholecystic fluid, and a positive sonographic Murphy sign. Treatment involves intravenous antibiotics and cholecystectomy, typically performed laparoscopically either early or delayed depending on clinical severity.

\medterm{Chronic Granulomatous Disease}-- See immunology.

\medterm{Chronic Inflammation}-- A prolonged inflammatory response characterized by simultaneous tissue destruction and attempts at repair, lasting weeks to years. Unlike acute inflammation, it is predominantly mononuclear, involving macrophages, lymphocytes, and plasma cells rather than neutrophils. The process is driven by persistent infections such as tuberculosis, autoimmune reactions as in rheumatoid arthritis, prolonged exposure to toxic agents including silica, or non-degradable foreign bodies. Chronic inflammation is characterized by mononuclear cell infiltration, tissue destruction, and fibrosis. Granuloma formation represents a specialized pattern seen in tuberculosis, sarcoidosis, and foreign body reactions. Chronic inflammation underlies many common diseases including rheumatoid arthritis and atherosclerosis.

\medterm{Cirrhosis}-- The end stage of chronic liver disease, characterized by diffuse fibrosis with distortion of the hepatic architecture and formation of regenerative nodules. The most common causes include chronic hepatitis C infection, alcoholic liver disease, non-alcoholic steatohepatitis, chronic hepatitis B infection, and autoimmune hepatitis. The pathogenesis involves repeated cycles of hepatocyte injury, inflammation, and repair, with activation of hepatic stellate cells that produce excessive extracellular matrix. The resulting fibrosis disrupts normal blood flow and hepatocyte function, causing portal hypertension and hepatic synthetic dysfunction. Clinical features include jaundice, ascites, variceal bleeding, hepatic encephalopathy, and coagulopathy. Cirrhosis is irreversible and increases the risk of hepatocellular carcinoma.

\medterm{Coagulation Cascade}-- The coagulation cascade is a series of enzymatic reactions leading to fibrin clot formation, involving serine protease activation cascades. The intrinsic pathway is initiated by contact activation when factor XII binds exposed collagen, sequentially activating factors XI, IX, and the cofactor VIII. The extrinsic pathway is initiated when tissue factor released from damaged cells binds factor VII, forming the TF-VIIa complex that activates factor X. Both pathways converge at factor X activation, forming the prothrombinase complex that converts prothrombin to thrombin. Thrombin cleaves fibrinogen to fibrin, which polymerises to form a stable clot. The cascade is regulated by natural anticoagulants including antithrombin III, protein C, and protein S. Laboratory assessment uses prothrombin time and activated partial thromboplastin time.

\medterm{Coagulative Necrosis}-- The most common form of necrosis, typically resulting from ischemic injury or infarction in most solid organs except the brain. The fundamental process involves denaturation of both structural proteins and enzymatic proteins, which prevents the immediate liquefaction of the dead tissue. The affected tissue maintains its basic structural outline for several days, a phenomenon known as structural preservation that allows histological identification of the original organ architecture. Over time, the necrotic tissue is removed by phagocytes and replaced by granulation tissue, which matures into a fibrous scar. Coagulative necrosis is the typical pattern in myocardial infarction, renal infarction, and splenic infarction, where the organ outline is preserved histologically for days to weeks.

\medterm{Colorectal Adenocarcinoma}-- The third most common cancer worldwide and the second leading cause of cancer death, arising from the colonic epithelium through the adenoma-carcinoma sequence over 10 to 15 years. The pathogenesis involves sequential accumulation of genetic mutations beginning with APC loss, followed by KRAS activation, TP53 loss, and other alterations. Approximately 15 percent of cases are associated with hereditary syndromes including Lynch syndrome with mismatch repair gene mutations, and familial adenomatous polyposis with APC mutation. Sporadic cases account for the majority. Screening via colonoscopy with polypectomy reduces both incidence and mortality. Treatment includes surgical resection, chemotherapy, and targeted therapies including anti-EGFR and anti-VEGF agents based on tumor molecular profile.

\medterm{Condition}-- A general term for a state of bodily health or, in medical usage, a specific deviation from normal in structure or function that may or may not constitute a disease, including complaints lacking a discrete pathologic basis and overt diagnoses. It is used when a precise disease cannot yet be identified and distinguished from a sign, symptom, or syndrome.

\medterm{Congestion}-- A passive vascular phenomenon of engorgement of a tissue or organ with blood from impaired venous or lymphatic drainage. Chronic venous congestion produces stasis, hypoxia, dilated vessels, edema, hemorrhage, and hemosiderin-laden macrophages, exemplified by the lungs (heart failure cells) and liver (nutmeg liver). Causes include heart failure, portal hypertension, and venous obstruction; longstanding congestion progresses to fibrosis.

\medterm{Cyanosis}-- A bluish discoloration of the skin and mucous membranes resulting from excessive amounts of deoxygenated hemoglobin in the capillary blood. It becomes clinically apparent when the concentration of deoxygenated hemoglobin exceeds 5 grams per deciliter, which corresponds to an oxygen saturation of approximately 85 percent. Central cyanosis affects the entire body and results from arterial hypoxemia, as seen in congenital heart disease with right-to-left shunting, severe pulmonary disease, and severe pneumonia. Peripheral cyanosis affects the extremities and results from reduced blood flow with increased oxygen extraction, as seen in cold exposure, shock, and peripheral vascular disease. Differential cyanosis, affecting only the lower body, can indicate specific congenital heart defects such as patent ductus arteriosus with Eisenmenger syndrome.

\medterm{Cyst}-- An abnormal closed cavity or sac lined by epithelium, containing liquid, semisolid, or gaseous material, and distinguished from an abscess (which contains pus) and a pseudocyst (which lacks an epithelial lining). Cysts occur in virtually every organ and tissue and are classified by pathogenesis: retention cysts (from obstruction of a duct, e.g., sebaceous cyst, ovarian follicular cyst, breast cyst), developmental cysts (arising from embryonic remnants, e.g., branchial cleft cyst, thyroglossal duct cyst, bronchogenic cyst), parasitic cysts (e.g., hydatid cyst from Echinococcus), and neoplastic cysts (e.g., cystic teratoma). Diagnosis is by imaging. Treatment depends on the type, size, location, and symptoms, ranging from observation to complete surgical excision.

\medterm{Cytology}-- The microscopic study of individual cells, distinct from histology in that tissue architecture is not examined, used for screening and diagnosis of neoplastic and infectious disease. Specimens are exfoliative (cervical smears, sputum) or aspirated (fine needle aspiration, body fluids), interpreted with Papanicolaou or special stains. Applications include the Papanicolaou test, sputum for lung cancer, and urine for urothelial carcinoma.

\medterm{Cytopathology}-- The branch of pathology using cytologic examination of cells to diagnose disease, particularly neoplastic and infectious processes, built on standardized reporting systems that guide management. Examples include the Bethesda System for thyroid fine needle aspiration and cervical cytology, and the Paris and Yokohama systems for urinary and pancreaticobiliary cytology. It overlaps with surgical pathology through immunocytochemistry and molecular testing.

\medterm{Dysplasia}-- Disordered cellular growth and maturation characterized by structural abnormalities within a tissue, representing a pre-neoplastic condition that may progress to invasive carcinoma. At the cellular level, dysplasia shows nuclear enlargement, hyperchromasia with clumped chromatin, irregular nuclear contours, increased nuclear-to-cytoplasmic ratio, and prominent nucleoli. Architecturally, there is loss of normal tissue organization, with cells failing to mature properly as they migrate from the basal layer to the surface. Dysplasia is graded as mild, moderate, or severe based on the extent of cytologic and architectural abnormalities. Severe dysplasia may progress to carcinoma in situ, where abnormal cells involve the full epithelial thickness but have not breached the basement membrane.

\medterm{Dystrophic Calcification}-- The deposition of calcium phosphate salts in damaged, necrotic, or degenerating tissues despite normal serum calcium levels. It occurs in areas of tissue injury, inflammation, or chronic degeneration where the local tissue environment favors calcium precipitation. The mechanism involves release of calcium-binding phospholipids and matrix vesicles from damaged cell membranes, which serve as nucleation sites for hydroxyapatite crystal formation. Common sites include atherosclerotic plaques, heart valves in degenerative aortic stenosis, and areas of previous trauma or infection. Dystrophic calcification can impair organ function, for example by restricting valve leaflet motion. It is distinct from metastatic calcification, which occurs in normal tissues in the setting of hypercalcaemia.

\medterm{Ecchymosis}-- An ecchymosis is a flat, purple or blue discoloration of the skin larger than 10 millimeters caused by extravasation of blood from damaged blood vessels into the subcutaneous tissue. Unlike petechiae, which result from platelet disorders, ecchymoses may result from coagulopathies, trauma, or vascular fragility. They are commonly seen in the elderly due to skin thinning and increased vascular fragility, as well as in patients taking anticoagulant or antiplatelet medications. The color of an ecchymosis evolves through characteristic changes over days to weeks: initially red-purple from deoxygenated hemoglobin, then blue-black, green from biliverdin, and finally yellow-brown from haemosiderin before complete resolution. The color progression helps estimate the age of the bruise.

\medterm{Edema}-- See general.

\medterm{Embolism}-- An embolism is the occlusion of a blood vessel by a detached mass that has traveled through the bloodstream from its site of origin. The most common form is thromboembolism, where a fragment of a thrombus dislodges and travels to a distant site through the venous or arterial circulation. Other types of emboli include fat emboli from bone fractures, amniotic fluid emboli during pregnancy, air emboli from trauma or surgery, tumor emboli from malignancy, foreign body emboli from injected materials, and cholesterol emboli from atherosclerotic plaque rupture. Clinical consequences depend on the size, composition, and location of the embolus. Pulmonary embolism is the most common life-threatening form. Treatment includes anticoagulation, thrombolysis, or embolectomy in selected cases.

\medterm{Exudate}-- An exudate is an inflammatory fluid with high protein content and abundant cells, resulting from increased vascular permeability during inflammation. Exudates are typically caused by bacterial infections, autoimmune diseases, viral infections, or malignancy. The fluid is often cloudy or turbid due to the presence of inflammatory cells, particularly neutrophils in bacterial infections. Exudates have high specific gravity above 1.020 and contain more than 300 cells per microliter. They are rich in fibrin and inflammatory mediators and are distinguished from transudates by Light's criteria: pleural fluid protein-to-serum protein ratio greater than 0.5, pleural fluid LDH-to-serum LDH ratio greater than 0.6, or pleural fluid LDH greater than two-thirds the upper limit of normal for serum.

\medterm{Fat Embolism}-- Fat embolism occurs when fat globules enter the bloodstream, typically following long bone fractures, orthopedic surgery, or severe trauma involving adipose tissue. The fat emboli travel through the venous circulation to the lungs, where they cause mechanical obstruction of pulmonary capillaries and chemical inflammation as lipases break down neutral fat into toxic free fatty acids. The pulmonary manifestations range from mild hypoxia to acute respiratory distress syndrome with severe hypoxemia and acute respiratory distress syndrome. Systemic embolisation causes petechial rash, neurological symptoms including confusion and focal deficits, and retinal changes. Diagnosis is clinical, supported by imaging. Treatment is supportive with oxygenation, ventilation, and hemodynamic support.

\medterm{Fat Necrosis}-- A specialized form of necrosis that occurs in adipose tissue, most commonly following acute pancreatitis when activated pancreatic lipases escape into the surrounding tissues and digest triglycerides in mesenteric or subcutaneous fat cells. The released fatty acids combine with calcium ions to form chalky-white deposits known as soap-stone lesions or saponification, which are visible on gross examination as firm, white, opaque areas. Histologically, fat necrosis shows shadowy outlines of necrotic adipocytes with basophilic calcium deposits. Over time, the necrotic fat is replaced by fibrosis and chronic inflammation. In the breast, fat necrosis can mimic carcinoma clinically and radiographically, often requiring biopsy for definitive diagnosis.

\medterm{Fever}-- See emergency_medicine.

\medterm{Fibrinoid Necrosis}-- A specialized pattern of vascular necrosis characterized by the deposition of immune complexes, fibrin, and plasma proteins within the walls of blood vessels, creating an amorphous, bright pink, smudgy eosinophilic appearance on hematoxylin and eosin staining. This pattern is characteristic of immune-mediated vascular injury, occurring in several important clinical conditions. In malignant hypertension, severely elevated blood pressure above 200/120 millimeters of mercury causes direct arterial wall damage, with plasma proteins leaking into the damaged vessel wall. In polyarteritis nodosa, immune complex deposition in small and medium arteries triggers fibrinoid necrosis. In systemic lupus erythematosus, similar immune complex-mediated vascular damage occurs. Affected vessels may thrombose, leading to tissue ischemia or rupture.

\medterm{Fibrinous Pericarditis}-- Inflammation of the pericardium characterized by deposition of fibrin on the pericardial surfaces, giving them a rough, shaggy appearance described as bread-and-butter pericarditis. It is the most common form of acute pericarditis, typically occurring after myocardial infarction, uremia, systemic lupus erythematosus, radiation therapy, or post-cardiac surgery syndrome. The fibrinous exudate consists of fibrin strands deposited on the visceral and parietal pericardial surfaces. The fibrin may organize into fibrous adhesions if the underlying process persists, potentially obliterating the pericardial space and causing constrictive pericarditis. Clinical presentation includes sharp retrosternal chest pain that improves with sitting forward, pericardial friction rub on auscultation, and diffuse ST-segment elevation on ECG.

\medterm{Fibroadenoma}-- The most common benign tumor of the female breast, a proliferation of stromal and epithelial elements presenting as a firm, rubbery, smooth, mobile, painless mass, most often in women aged fifteen to thirty. The mass slips easily under the fingers (breast mouse), and ultrasonography shows a well-defined oval hypoechoic lesion. Histologically it comprises fibrous stroma with benign epithelial-lined ducts.

\medterm{Fibrosis}-- The excessive deposition of extracellular matrix components, primarily collagen types I and III, in response to chronic tissue injury. It represents a pathological wound healing response that leads to organ scarring and dysfunction. The process is driven by activated fibroblasts and myofibroblasts, which are stimulated by growth factors including TGF-beta, PDGF, and CTGF produced by inflammatory cells and injured epithelium. Fibrosis can occur in any organ, but is particularly problematic in the liver (cirrhosis), lungs (pulmonary fibrosis), heart (myocardial fibrosis), and kidneys (renal fibrosis). Unlike normal wound healing, fibrosis involves excessive extracellular matrix that distorts tissue architecture and impairs organ function. Fibrosis is often progressive and currently lacks effective treatments beyond addressing the underlying cause.

\medterm{Fine Needle Aspiration (FNA)}-- A percutaneous sampling technique in which a 22 to 25-gauge needle, often ultrasound- or palpation-guided, aspirates cellular material from a palpable or imaging-detected lesion for cytologic examination, with minimal trauma and lower morbidity than core or open biopsy. Widely used for thyroid nodules, breast lumps, lymph nodes, salivary glands, and soft tissue lesions; rapid on-site evaluation improves adequacy.

\medterm{Fistula}-- A fistula is an abnormal epithelial-lined connection between two body surfaces or organs that normally do not communicate. Fistulas may be congenital, such as tracheoesophageal fistula, or acquired, resulting from inflammation, infection, surgery, trauma, or malignancy. Common acquired types include enterocutaneous fistulas between the intestine and skin, seen after surgery or in Crohn's disease; perianal fistulas following anorectal abscesses; and arteriovenous fistulas between arteries and veins from traumatic injury or congenital malformations. Fistulas may cause continuous leakage of contents, infection, electrolyte disturbances, and malnutrition. Diagnosis involves imaging including fistulography, CT, or MRI. Treatment depends on the site and cause, ranging from medical therapy to surgical closure.

\medterm{FNA}-- See Fine Needle Aspiration (FNA).

\medterm{Free Radical Injury}-- Free radicals are highly reactive molecules with unpaired electrons that can damage cellular components through oxidative stress. Reactive oxygen species include the superoxide anion, hydrogen peroxide, and the hydroxyl radical generated through the Fenton reaction. Sources include normal mitochondrial oxidative metabolism, the NADPH oxidase respiratory burst in neutrophils, cytochrome P450 reactions, and external sources including ionizing radiation and ultraviolet light. Free radical damage to lipids initiates lipid peroxidation of polyunsaturated fatty acids in cell membranes, disrupting membrane integrity and function. Protein damage causes enzyme inactivation and receptor dysfunction. DNA damage leads to strand breaks and base modifications. Antioxidant defences include superoxide dismutase, catalase, glutathione peroxidase, and vitamins C and E.

\medterm{Frozen Section}-- An intraoperative pathology consultation in which a fresh tissue specimen is rapidly frozen, sectioned on a cryostat, and stained with hematoxylin and eosin within 10 to 20 minutes to guide immediate surgical decision-making. It is used to establish benign versus malignant, assess margins, and identify lymph node metastases. Freezing artifacts make it less accurate than permanent paraffin histology.

\medterm{Gangrene}-- Death of tissue due to loss of blood supply, typically affecting the extremities and bowel, and represents a form of coagulative necrosis. Dry gangrene occurs in tissues with gradual arterial occlusion, where the dead tissue becomes dry, shrunken, and dark brown or black. It is commonly seen in the toes and feet of patients with diabetes mellitus and peripheral arterial disease. The black color results from hemoglobin degradation and iron sulfide deposition in the necrotic tissue. A clear line of demarcation typically separates viable from non-viable tissue. Wet gangrene results from bacterial superinfection of ischemic tissue, producing a swollen, foul-smelling, liquefied mass with systemic toxicity and a less distinct line of demarcation. Gas gangrene is a specific form caused by Clostridium species, particularly C. perfringens, producing gas within tissues and representing a surgical emergency requiring urgent debridement and antibiotics.

\medterm{Gas Gangrene}-- A life-threatening bacterial infection characterized by the presence of gas in tissues (subcutaneous emphysema) and rapid myonecrosis. Caused by Clostridium perfringens (type A, 80\% of cases), C. novyi, C. septicum, and C. histolyticum--Gram-positive, spore-forming, obligate anaerobic bacilli that produce exotoxins (alpha-toxin: phospholipase C, destroys cell membranes; theta-toxin: cytolysin; collagenase, hyaluronidase, DNase). Predisposing factors: contaminated wounds (soil, foreign bodies), trauma (crush injuries, deep lacerations), devascularised tissue, immunosuppression, and abdominal surgery. Incubation period: 6 hours to 6 days. Clinical features: sudden onset of severe pain at wound site, discoloured (bronze or purplish) skin, bullae filled with serosanguinous fluid, crepitus (gas in tissue, palpable), foul-smelling discharge, tachycardia (out of proportion to fever), hypotension, renal failure, hemolytic anemia, sepsis, and rapidly progressive necrosis. Diagnosis: clinical (crepitus, gas on X-ray/CT, Gram stain of exudate showing Gram-positive rods with absent white cells), culture, and PCR. Treatment: emergent surgical debridement (fasciotomy and excision of all necrotic tissue, often requiring amputation), high-dose intravenous penicillin G plus clindamycin (metronidazole as alternative), hyperbaric oxygen therapy (adjunctive, suppresses toxin production by raising tissue oxygen tension). Mortality: 20-30\% overall, >50\% with trunk involvement or delay in treatment.

\medterm{Glomerulonephritis}-- See nephrology.

\medterm{Granuloma}-- See immunology.

\medterm{Granulomatous Inflammation}-- A distinctive pattern of chronic inflammation characterized by organized collections of activated macrophages that have undergone morphological transformation into epithelioid cells, often surrounded by a rim of lymphocytes and occasionally fibroblasts. Epithelioid cells have abundant eosinophilic cytoplasm and cell borders that blend together, resembling epithelial cells. Multinucleated giant cells, formed by fusion of macrophages under the influence of cytokines including interferon-gamma, form around persistent stimuli such as Mycobacterium tuberculosis, fungi, foreign bodies, and in sarcoidosis. The granuloma walls off the inciting agent and contains the infection, but the inflammatory response also causes tissue destruction and fibrosis in surrounding tissues.

\medterm{Gross Examination}-- The macroscopic, naked-eye inspection of a specimen in the pathology laboratory, performed before microscopic sectioning and the essential first step in pathologic diagnosis of resected tissues and organs. It records the specimen's origin, dimensions, weight, color, shape, and lesion characteristics, and determines how it is sampled for histology. Accurate gross description underpins microscopic interpretation, staging, and margin assessment.

\medterm{Hemorrhage}-- Extravasation of blood from the cardiovascular system (heart, arteries, veins, capillaries). Classification by severity: class 1 (<15\% blood loss, compensated, minimal symptoms), class 2 (15-30\%, mild tachycardia, decreased pulse pressure), class 3 (30-40\%, significant tachycardia, hypotension, pallor, oliguria, altered mental status--decompensated), class 4 (>40\%, severe hypotension, anuria, loss of consciousness--immediately life-threatening). Classification by location: external (visible from wound, or body orifice--epistaxis, haemoptysis, hematemesis, melaena, haematuria, vaginal) vs internal (concealed, into body cavities--intracranial, intrathoracic/ haemothorax, intra-abdominal/ haemoperitoneum, retroperitoneal, haemarthrosis, muscle haematoma). Causes: trauma (most common cause of preventable death from exsanguination), ruptured aneurysm (AAA, cerebral berry aneurysm, splenic artery), gastrointestinal bleeding (peptic ulcer, varices, diverticulosis, angiodysplasia, cancer), obstetric hemorrhage (postpartum hemorrhage--atomic uterus, placental abruption, placenta praevia), surgery, anticoagulation, coagulopathy (hemophilia, DIC, vitamin K deficiency). Management: ABC (airway, breathing, circulation), direct pressure for external bleeding, IV access (two large-bore cannulae), fluid resuscitation (crystalloids, blood products--massive transfusion protocol: 1:1:1 ratio of PRBCs:FFP:platelets), pressure control (tourniquet for life-threatening extremity hemorrhage, pelvic binder for pelvic fracture), reversal of anticoagulation (vitamin K, fresh frozen plasma, PCC--prothrombin complex concentrate, protamine for heparin, tranexamic acid), damage control surgery (abbreviated laparotomy, packing to control hemorrhage, angioembolisation, endovascular stenting).

\medterm{Hemosiderin}-- A golden-brown, granular, iron-containing pigment derived from the degradation of hemoglobin and stored within macrophages as an insoluble complex. It forms when red blood cells are phagocytosed and degraded, typically in conditions of hemorrhage, hemolysis, or excessive red blood cell destruction. The iron from hemoglobin is initially stored as soluble ferritin, but under conditions of iron overload, ferritin aggregates to form hemosiderin, which is more insoluble and visible on routine haematoxylin and eosin sections. The Prussian blue stain specifically identifies iron in haemosiderin deposits. Common sites include the lungs in chronic left heart failure (heart failure cells), the spleen and liver in chronic hemolytic anemia, and at sites of old hemorrhage.

\medterm{Histology}-- The microscopic study of the structure of tissues and cells. Histology involves preparing thin tissue sections (3-10 micrometres) by fixation (formalin, paraffin embedding), sectioning (microtome), staining (haematoxylin and eosin, H\&E), and examination under a light or electron microscope. Basic tissue types: (1) Epithelium: lining of surfaces and cavities, glandular tissue; classified by cell shape (squamous, cuboidal, columnar) and layering (simple, stratified, pseudostratified). (2) Connective tissue: supports and connects other tissues; includes bone, cartilage, adipose, blood, and loose/dense connective tissue; consists of cells (fibroblasts, osteocytes, chondrocytes, adipocytes, macrophages, mast cells) embedded in extracellular matrix (collagen, elastin, proteoglycans). (3) Muscle tissue: skeletal (striated, voluntary), cardiac (striated, involuntary, intercalated discs), smooth (non-striated, involuntary). (4) Nervous tissue: neurons (conduct electrical impulses) and glial cells (support, myelinate, nourish, defend). Special stains: Masson trichrome (collagen), periodic acid-Schiff (PAS, glycogen/mucin), silver stains (reticulin, nerve fibers), immunohistochemistry (antigen detection using labelled antibodies). Histopathology is the histological diagnosis of disease (cancer grading/staging, identification of inflammatory or infectious processes).

\medterm{Histopathology}-- The laboratory-based diagnosis of disease through microscopic examination of tissue sections from biopsy, excision, or autopsy, built on fixation, paraffin embedding, sectioning, and hematoxylin and eosin staining. Interpretation integrates architecture with morphology to identify inflammatory, infectious, metabolic, degenerative, and neoplastic processes and to assign tumor type, grade, and stage. Adjuncts include special stains, immunohistochemistry, and molecular assays.

\medterm{Huntington's Disease}-- An autosomal dominant neurodegenerative disorder caused by a CAG trinucleotide repeat expansion in the huntingtin gene on chromosome 4p16.3. The disease is characterized by the triad of progressive chorea, cognitive decline, and psychiatric manifestations. The pathogenesis involves accumulation of mutant huntingtin protein with expanded polyglutamine tract in neurons, particularly in the caudate nucleus and putamen, leading to transcriptional dysregulation, mitochondrial dysfunction, excitotoxicity, and oxidative stress. The disease typically manifests in mid-adulthood with progressive chorea, cognitive decline, and psychiatric symptoms. The number of CAG repeats correlates inversely with age of onset. There is no cure; treatment is supportive and symptom-focused, with tetrabenazine for chorea.

\medterm{Hyperaldosteronism}-- Hyperaldosteronism, also known as Conn syndrome when caused by an adrenal adenoma, is characterized by excessive aldosterone production causing hypertension, hypokalemia, and metabolic alkalosis. Primary hyperaldosteronism results from autonomous aldosterone production by the adrenal gland, most commonly from a unilateral aldosterone-producing adenoma or bilateral adrenal hyperplasia. The excessive aldosterone acts on the distal nephron, promoting sodium reabsorption and potassium and hydrogen ion secretion, leading to hypokalaemia and metabolic alkalosis. Hypertension results from sodium and water retention. Diagnosis involves measuring the plasma aldosterone-to-renin ratio, followed by confirmatory suppression testing. Treatment includes surgical resection for adenomas and mineralocorticoid receptor antagonists for bilateral hyperplasia.

\medterm{Hypersensitivity Type I}-- Type I hypersensitivity, also known as immediate hypersensitivity or allergic reactions, is mediated by IgE antibodies and mast cells, producing symptoms within minutes of allergen re-exposure. During sensitization, initial allergen exposure activates Th2 cells that secrete IL-4 and IL-13, promoting B cell class switching to IgE production. IgE antibodies bind with high affinity to Fc epsilon RI receptors on mast cells and basophils, priming these cells for re-exposure. Upon subsequent allergen exposure, IgE cross-linking on mast cells triggers degranulation, releasing histamine, tryptase, leukotrienes, and prostaglandins that cause vasodilation, bronchoconstriction, and increased vascular permeability. Clinical manifestations range from mild allergic rhinitis and urticaria to life-threatening anaphylaxis. Treatment includes antihistamines, corticosteroids, and epinephrine.

\medterm{Hypersensitivity Type II}-- Type II hypersensitivity, also known as cytotoxic hypersensitivity, is mediated by IgG or IgM antibodies directed against antigens on cell surfaces or the extracellular matrix, causing tissue damage through complement activation, antibody-dependent cellular cytotoxicity, and phagocytosis. When antibodies bind to cell-surface antigens, the classical complement pathway is activated, generating C3b for opsonization and the membrane attack complex for direct cell lysis. Antibody-coated cells are also killed by NK cells through antibody-dependent cellular cytotoxicity. Clinical examples include autoimmune hemolytic anemia, Goodpasture syndrome, myasthenia gravis, and immune thrombocytopenia. In drug-induced hemolytic anemia, drugs such as penicillin act as haptens, binding to red cell surfaces and triggering antibody-mediated destruction.

\medterm{Hypersensitivity Type III}-- Type III hypersensitivity is caused by deposition of antigen-antibody immune complexes in tissues and blood vessel walls, activating complement and recruiting neutrophils that release destructive enzymes and reactive oxygen species. This reaction occurs when there is a slight excess of antigen relative to antibody, forming small, soluble immune complexes that are not efficiently cleared by the reticuloendothelial system. The complexes deposit in vessel walls, particularly in the kidneys, joints, and skin, causing vasculitis, arthritis, and glomerulonephritis. Clinical examples include systemic lupus erythematosus (DNA-anti-DNA complexes), post-streptococcal glomerulonephritis, serum sickness, and the Arthus reaction. Chronic immune complex deposition leads to persistent inflammation and progressive tissue damage.

\medterm{Hypersensitivity Type IV}-- Type IV hypersensitivity, also known as delayed-type hypersensitivity, is mediated by T lymphocytes rather than antibodies and takes 48 to 72 hours to develop after antigen exposure. The reaction involves sensitized CD4 Th1 cells that, upon re-exposure to the antigen presented by dendritic cells or macrophages, release pro-inflammatory cytokines including interferon-gamma and TNF, which activate macrophages and cause tissue damage. Contact dermatitis is the most common clinical example, occurring after contact with poison ivy, nickel, and other sensitising chemicals. The tuberculin reaction (PPD test) is another classic example. Type IV hypersensitivity is also involved in allograft rejection and many autoimmune diseases. It is characterized by a delayed onset of 24 to 72 hours.

\medterm{Hypovolemic Shock}-- See emergency_medicine.

\medterm{Increased Vascular Permeability}-- A fundamental feature of acute inflammation that allows plasma proteins, fluid, and leukocytes to exit the vascular space and enter the extravascular tissue. The most common mechanism involves endothelial cell contraction caused by histamine, leukotrienes C4 and D4, and complement components C3a and C5a, creating intercellular gaps between endothelial cells that permit passage of protein-rich fluid. This response occurs predominantly in postcapillary venules and is immediate but short-lived, lasting 15 to 30 minutes. Other mechanisms include direct endothelial damage (burns, radiation, toxins) causing immediate and prolonged leakage, leukocyte-mediated damage via toxic oxygen species, and increased transcytosis through vesicular transport.

\medterm{Inflammatory Myopathies}-- A group of autoimmune disorders characterized by chronic muscle inflammation and progressive proximal muscle weakness. The three main types are dermatomyositis, polymyositis, and inclusion body myositis. Dermatomyositis is characterized by perifascicular atrophy of muscle fibers, perivascular inflammation, and characteristic skin findings including heliotrope rash on the eyelids, Gottron papules over the knuckles, and mechanic hands. It is associated with complement-mediated microangiopathy with membrane attack complex deposition in endomysial capillaries. Polymyositis shows endomysial inflammation with CD8 T cells invading non-necrotic muscle fibers. Inclusion body myositis, the most common type in patients over 50, features rimmed vacuoles and protein aggregates. Diagnosis involves EMG, muscle biopsy, and elevated creatine kinase.

\medterm{Inspection}-- The first step of the physical examination, involving systematic visual observation of the patient\'s overall appearance, posture, body habitus, gait, skin color, nutritional status, respiratory effort, and specific anatomical regions. Requires adequate lighting and proper patient exposure. Provides immediate diagnostic clues prior to palpation, percussion, or auscultation.

\medterm{Interstitial Nephritis}-- Inflammation of the renal interstitium and tubules, which can be acute or chronic. Acute interstitial nephritis is most commonly caused by medications, particularly NSAIDs, antibiotics including penicillins and cephalosporins, and proton pump inhibitors. The pathogenesis involves immune-mediated hypersensitivity reaction to the offending drug, with T-cell and eosinophil infiltration of the renal interstitium. Clinical features include acute kidney injury occurring days to weeks after starting the offending medication, with fever, rash, eosinophilia, and sterile pyuria. Renal biopsy shows interstitial inflammation with eosinophils. Treatment involves discontinuation of the causative drug and corticosteroids in severe cases. Most patients recover renal function fully.

\medterm{Ischemia-Reperfusion Injury}-- The paradoxical tissue damage that occurs when blood flow is restored to previously ischemic tissue, a phenomenon with major clinical significance in myocardial infarction treatment, stroke management, organ transplantation, and shock resuscitation. During ischemia, ATP depletion causes failure of the sodium-potassium pump leading to cellular swelling, calcium accumulation, activation of phospholipases and proteases, and anaerobic metabolism with lactic acid accumulation. Upon reperfusion, the sudden reintroduction of oxygen triggers a burst of reactive oxygen species from mitochondria and neutrophils, causing lipid peroxidation, protein damage, and DNA injury. Inflammatory mediators are released, recruiting more neutrophils and amplifying the damage. Therapeutic strategies focus on minimising this injury through controlled reperfusion.

\medterm{Lesion}-- A general term for any abnormal change in tissue structure due to disease or injury. Lesions can occur in any organ or tissue and can be classified by: (1) morphology: macule, papule, plaque, nodule, vesicle, bulla, pustule, ulcer, erosion, fissure, tumor; (2) location: skin, organ-specific, or systemic; (3) etiology: traumatic, infectious, inflammatory, neoplastic (benign or malignant), degenerative, congenital, vascular, metabolic; (4) imaging characteristics: radiolucent, radiopaque, contrast-enhancing, cystic, solid. A lesion is a descriptive finding, not a diagnosis; further characterisation (biopsy, imaging) is needed to determine the specific pathology.

\medterm{Lipofuscin}-- A yellow-brown granular pigment composed of oxidized lipids and cross-linked proteins that accumulates progressively within cells with age, particularly in post-mitotic cells such as neurons, cardiac myocytes, and hepatocytes. It represents the residues of lysosomal digestion of autophagocytosed organelles and damaged cellular components that cannot be further degraded or exocytosed. Lipofuscin accumulation is a hallmark of cellular aging and is found in increased amounts in tissues with high oxidative stress. It is often called the 'wear and tear' pigment of ageing. Lipofuscin granules are autofluorescent and stain with periodic acid-Schiff and Sudan black. Its accumulation does not directly harm cells but serves as a marker of cumulative oxidative damage and cellular senescence.

\medterm{Liquefactive Necrosis}-- Characterized by the complete enzymatic digestion of dead cells, resulting in a liquid, viscous mass that eventually liquefies and forms a cystic cavity. This form of necrosis occurs most commonly in focal bacterial or fungal infections, where the inflammatory response releases abundant hydrolytic enzymes from neutrophils and macrophages that digest both dead and viable tissue. In the central nervous system, liquefactive necrosis follows ischemic stroke because the brain has abundant hydrolytic enzymes and limited connective tissue, so necrotic tissue liquefies rapidly, forming a cystic cavity. In bacterial infections, accumulated neutrophils release proteolytic enzymes that digest dead tissue, producing pus. The resulting cavity may persist as a cyst or be replaced by fibrous tissue during healing.

\medterm{Liver Failure}-- A life-threatening condition resulting from severe impairment of hepatic synthetic, metabolic, and detoxification functions. Acute liver failure occurs in previously healthy individuals, most commonly from acetaminophen toxicity accounting for approximately 50 percent of cases in the United States, viral hepatitis, drug-induced liver injury, or autoimmune hepatitis. The pathogenesis involves massive hepatocyte necrosis leading to loss of synthetic function with coagulopathy, impaired bilirubin clearance with jaundice, and hepatic encephalopathy from failure to detoxify ammonia and other neurotoxins. Chronic liver failure develops more slowly on a background of cirrhosis, with portal hypertension and synthetic dysfunction. Liver transplantation is the definitive treatment for end-stage liver disease.

\medterm{Lupus Erythematosus}-- Systemic lupus erythematosus is a chronic autoimmune disease characterized by production of autoantibodies against nuclear antigens, leading to immune complex-mediated inflammation and tissue damage in multiple organ systems. The disease predominantly affects women of childbearing age, with a striking female predominance of 9 to 1. The pathogenesis involves loss of self-tolerance to nuclear antigens, with formation of immune complexes that deposit in tissues and activate complement, causing inflammation in the kidneys, skin, joints, and serosal surfaces. The malar butterfly rash is characteristic. Key laboratory findings include positive ANA, anti-dsDNA, and anti-Smith antibodies. Disease flares and remissions are typical. Treatment includes NSAIDs, corticosteroids, hydroxychloroquine, and immunosuppressive agents for severe organ involvement.

\medterm{Malignant Neoplasm}-- Malignant neoplasms, commonly called cancers, are uncontrolled proliferations of cells that invade surrounding tissues and have the capacity to metastasize to distant sites through lymphatic, hematogenous, or transcoelomic routes. They exhibit infiltrative growth, penetrating basement membranes and spreading along tissue planes, nerve sheaths, and vascular structures. Malignant cells show varying degrees of anaplasia, increased and often atypical mitotic activity, and loss of normal tissue architecture capable of invasion and metastasis. Malignant tumors grow rapidly, have poorly defined margins, and can recur after resection. They cause morbidity through local tissue destruction, obstruction, hemorrhage, and systemic effects including cachexia, paraneoplastic syndromes, and metastatic spread.

\medterm{Mediators of Inflammation}-- Inflammatory mediators are soluble molecules that initiate, amplify, or regulate the inflammatory response. They include lipid mediators derived from arachidonic acid such as prostaglandins and leukotrienes, vasoactive amines like histamine and serotonin, cytokines including TNF, IL-1, and IL-6, chemokines that direct leukocyte migration, complement components, and kinins like bradykinin. These mediators are produced by activated macrophages, mast cells, endothelial cells, and platelets. Many mediators have multiple overlapping functions and are tightly regulated to prevent excessive tissue damage. Understanding their roles has enabled targeted therapies including NSAIDs (inhibiting prostaglandins), leukotriene inhibitors, and biologic agents blocking TNF-alpha, IL-1, and IL-6 in chronic inflammatory diseases such as rheumatoid arthritis and inflammatory bowel disease.

\medterm{Metaplasia}-- A reversible adaptive change in which one differentiated adult cell type is replaced by another, representing the body's attempt to better withstand a chronic stress or irritation. The process involves reprogramming of tissue stem cells to differentiate along a different lineage that is more resilient to the new environment. The most common example is squamous metaplasia in the respiratory tract of chronic smokers, where the normal ciliated pseudostratified columnar epithelium is replaced by non-ciliated stratified squamous epithelium that is more resistant to smoke damage. Another key example is Barrett's esophagus, where the normal squamous epithelium of the lower esophagus is replaced by intestinal-type columnar epithelium due to chronic gastro-esophageal reflux, increasing the risk of esophageal adenocarcinoma.

\medterm{Metastatic Calcification}-- The deposition of calcium salts in otherwise normal tissues due to systemic hypercalcemia, where elevated serum calcium levels exceed the tissue ability to maintain calcium homeostasis. Causes include hyperparathyroidism with excess PTH causing bone resorption, vitamin D intoxication increasing intestinal calcium absorption, milk-alkali syndrome, bone destruction from malignancy or Paget's disease, sarcoidosis through extrarenal production of 1,25-dihydroxyvitamin D, and renal failure with phosphate retention. Metastatic calcification primarily affects tissues that excrete acid or generate alkaline environments, including the gastric mucosa, kidneys, lungs, and blood vessels. The deposits are intracellular and extracellular and may impair organ function over time.

\medterm{Microscopic Examination}-- The pathologic evaluation of tissue or cells under the light or electron microscope, complementing the gross examination to form the complete morphologic diagnosis. Routine light microscopy uses hematoxylin and eosin sections to assess architecture, cellular morphology, mitotic activity, and relation to margins. Special stains (periodic acid-Schiff, Masson trichrome, reticulin, Grocott methenamine silver) highlight specific components; immunohistochemistry adds further detail.

\medterm{Mural Thrombus}-- A mural thrombus is a blood clot that forms on the wall of a heart chamber or large blood vessel, typically in areas of damaged endothelium, stasis, or turbulent flow. In the heart, mural thrombi commonly form in the left ventricle following myocardial infarction, where the akinetic or dyskinetic segment promotes stasis of blood, or in the left atrium in patients with atrial fibrillation or mitral stenosis where blood stasis favors clot formation. The thrombus is attached to the wall by a broad base. In large vessels, mural thrombi form over atherosclerotic plaques or aneurysms. They can embolise to distant sites, causing infarction. Treatment includes anticoagulation to prevent propagation and embolisation. Over time, mural thrombi may organize via fibroblast ingrowth and become incorporated into the vessel wall or recanalise.

\medterm{Myocardial Infarction}-- The irreversible necrosis of myocardial tissue resulting from prolonged ischemia, most commonly due to atherosclerotic plaque rupture with superimposed thrombosis in a coronary artery. The process begins with platelet adhesion to the disrupted plaque, followed by activation of the coagulation cascade and formation of an occlusive thrombus. The resulting complete cessation of blood flow causes coagulative necrosis of the supplied myocardium within 20 to 40 minutes. The infarct evolves through stages: coagulation necrosis in the first 24 hours, neutrophil infiltration at days 1 to 3, granulation tissue formation at days 7 to 14, and eventual fibrous scar formation over 4 to 6 weeks. Diagnosis relies on ECG changes, elevated cardiac troponin, and imaging. Treatment includes reperfusion, antiplatelet therapy, and risk factor modification.

\medterm{Necroptosis}-- A regulated, programmed form of necrosis that shares features of both classical necrosis and apoptosis, serving as a backup cell death mechanism when apoptosis is inhibited. It is particularly important during viral infections, as many viruses encode proteins that block caspase-mediated apoptosis to ensure their own replication. The process is mediated by receptor-interacting protein kinase 1 and 3, known as RIPK1 and RIPK3, which form a complex called the necrosome. When TNF binds its receptor, RIPK1 and RIPK3 are recruited and activated. RIPK3 phosphorylates MLKL, which oligomerises and inserts into the plasma membrane, forming pores that cause cell swelling and membrane rupture, releasing damage-associated molecular patterns that trigger inflammation. Necroptosis is implicated in ischemia-reperfusion injury and neurodegenerative diseases.

\medterm{Necrosis}-- A pathological, unregulated form of cell death resulting from severe, irreversible injury that invariably provokes an intense inflammatory response. Unlike apoptosis, necrosis occurs when cells are exposed to extreme stresses such as profound hypoxia, toxins, trauma, burns, or extreme temperatures that overwhelm cellular repair mechanisms. The hallmark features include cellular swelling due to failure of the sodium-potassium pump, membrane disruption, organelle breakdown, and leakage of cytoplasmic contents into the extracellular space, triggering an inflammatory response. The morphologic patterns include coagulative (most common, in solid organs), liquefactive (brain and abscesses), caseous (tuberculosis), fat (adipose tissue), and fibrinoid (blood vessels).

\medterm{Neurodegeneration}-- The progressive loss of structure and function of neurons, culminating in their death, and represents the underlying mechanism of many devastating neurological disorders. It encompasses a heterogeneous group of diseases including Alzheimer's disease, Parkinson's disease, Huntington's disease, amyotrophic lateral sclerosis, and frontotemporal dementia. The mechanisms of neuronal death are multifactorial and include protein aggregation and misfolding, mitochondrial dysfunction, and oxidative stress. Common themes across neurodegenerative diseases include aberrant protein aggregation, mitochondrial impairment, excitotoxicity, neuroinflammation, and defective axonal transport. Despite intensive research, most neurodegenerative diseases remain incurable, with current treatments focused on symptom management.

\medterm{Neurotransmitter}-- Neurotransmitters are endogenous chemical messengers that transmit signals across synapses from one neuron to another or to target cells such as muscle or gland cells. They are synthesized in the presynaptic neuron from precursor molecules through specific enzymatic pathways, stored in synaptic vesicles, and released upon arrival of an action potential and calcium influx. Major neurotransmitters include acetylcholine synthesized from choline and acetyl-CoA, dopamine synthesized from tyrosine through tyrosine hydroxylase, and serotonin from tryptophan through tryptophan hydroxylase. After release, neurotransmitters bind to specific receptors on the postsynaptic membrane, either ionotropic receptors that directly open ion channels or metabotropic G-protein-coupled receptors. Signalling is terminated by reuptake, enzymatic degradation, or diffusion.

\medterm{Neutrophil Recruitment}-- A multistep process essential for acute inflammation and host defense, involving sequential adhesive interactions between neutrophils and endothelial cells. The process begins with the rolling phase, where neutrophils loosely attach to endothelial cells through selectin-mediated interactions. Activated endothelial cells express P-selectin stored in Weibel-Palade bodies and newly synthesized E-selectin, which bind to sialyl-Lewis X carbohydrates on neutrophils. The next step is firm adhesion, mediated by integrins including LFA-1 and Mac-1 on neutrophils binding to ICAM-1 on activated endothelial cells. Finally, neutrophils transmigrate through the endothelium via the paracellular route between endothelial cells (diapedesis) and migrate along chemotactic gradients to reach the site of injury.

\medterm{Obstructive Shock}-- A form of shock caused by physical obstruction of blood flow outside the heart, reducing cardiac output despite normal cardiac function. Common causes include cardiac tamponade, where fluid accumulation in the pericardial space compresses the heart and prevents adequate diastolic filling; tension pneumothorax, where air in the pleural space causes mediastinal shift and compression of the great veins reducing venous return; and massive pulmonary embolism, where thrombus in the pulmonary arteries obstructing right ventricular outflow. In all cases, cardiac output falls and compensatory vasoconstriction may be insufficient. Treatment involves relieving the obstruction: pericardiocentesis for tamponade, needle decompression for tension pneumothorax, and thrombolysis or embolectomy for massive pulmonary embolism.

\medterm{Obstructive Uropathy}-- Urinary tract obstruction at any level from the renal pelvis to the urethra, leading to impaired urine flow and potential renal damage. Causes include nephrolithiasis, benign prostatic hyperplasia, ureteral strictures, tumors of the bladder, prostate, or cervix, and external compression from lymphadenopathy or retroperitoneal fibrosis. The pathogenesis involves increased intratubular pressure proximal to the obstruction, leading to decreased glomerular filtration, tubular dysfunction, and ultimately irreversible renal parenchymal damage (obstructive nephropathy). Acute complete obstruction causes severe flank pain, anuria, and rapidly progressive renal failure. Chronic partial obstruction leads to hydronephrosis and gradual loss of renal function. Treatment involves relieving the obstruction.

\medterm{Ochronosis}-- A blue-black or brownish pigmentation of connective tissues (cartilage, skin, sclerae, tendons) caused by accumulation of homogentisic acid and its oxidation products. Ochronosis is the characteristic feature of alkaptonuria (also spelled alcaptonuria), a rare autosomal recessive disorder of tyrosine metabolism due to deficiency of homogentisate 1,2-dioxygenase (HGD gene). Homogentisic acid accumulates, polymerises, and deposits in collagen-rich tissues. Clinical: urine turns dark on standing (alkalinisation or oxidation—classic sign), ochronotic pigmentation (ear cartilage—blue-black, sclerae—triangular brown patches, skin over nose and knuckles, costochondral junctions), and ochronotic arthropathy (progressive osteoarthritis, especially of the spine and large joints—lumbar kyphosis, disc degeneration with calcification on X-ray, joint space narrowing). Other complications: aortic valve stenosis, kidney stones, and prostate stones. Onset of arthropathy in the 4th--5th decade. Treatment: nitisinone (inhibits 4-hydroxyphenylpyruvate dioxygenase upstream, reduces homogentisic acid production—licensed for alkaptonuria), low-protein diet, vitamin C (may delay polymerisation), and joint replacement for advanced arthropathy.

\medterm{Oncogenes}-- Mutated, amplified, or overexpressed versions of normal cellular genes called proto-oncogenes that encode proteins involved in cell growth, proliferation, and survival. Proto-oncogenes encode growth factors, growth factor receptors, signal transduction molecules, and transcription factors that normally regulate controlled cell growth. When activated by mutation, gene amplification, chromosomal translocation, or viral insertion, they become constitutively active, driving uncontrolled cell proliferation and survival. Common oncogenes include RAS (mutated in many cancers), MYC (transcription factor), HER2/ERBB2 (amplified in breast cancer), and BCR-ABL (translocation causing chronic myeloid leukemia). Targeted therapies against specific oncogenic drivers, such as tyrosine kinase inhibitors, have revolutionised cancer treatment.

\medterm{Oncosis}-- A form of cell death characterized by cellular swelling, derived from the Greek word onkos meaning swelling, representing a pre-lethal state that precedes necrosis. It is typically caused by ischemic injury leading to ATP depletion, which impairs the sodium-potassium pump and causes osmotic water influx. Unlike apoptosis, oncosis involves progressive cellular and organellar swelling, increased membrane permeability, cellular fragmentation, and eventual membrane rupture. Histologically, oncotic cells appear swollen with pale, vacuolated cytoplasm and clumped chromatin. The term is often used interchangeably with necrosis in the context of ischemic injury. Oncosis is distinct from apoptosis, which involves cell shrinkage rather than swelling, and is considered an unregulated, energy-independent form of cell death.

\medterm{p53 Tumor Suppressor}-- The p53 protein, encoded by the TP53 gene on chromosome 17p, is a transcription factor and the most commonly mutated gene in human cancer, with mutations found in over 50 percent of all malignancies. Known as the guardian of the genome, p53 is a short-lived protein normally kept at low levels by MDM2-mediated ubiquitination and proteasomal degradation. In response to DNA damage, oncogene activation, hypoxia, or nutrient deprivation, p53 is stabilized through phosphorylation and acetylation that increase its half-life and transcriptional activity. Activated p53 transactivates target genes involved in cell cycle arrest (p21), DNA repair (GADD45), apoptosis (BAX, PUMA), and senescence. Mutant p53 not only loses tumor suppressor function but often gains oncogenic properties. Germline TP53 mutations cause Li-Fraumeni syndrome.

\medterm{Palpation}-- The second step of the physical examination, utilizing the sense of touch to evaluate body structures for size, shape, consistency, mobility, tenderness, temperature, and vibration. Technique: light palpation (1-2 cm depth) assesses superficial tenderness and muscle guarding; deep palpation (4-5 cm depth) delineates abdominal organomegaly and masses. Palmar surfaces of fingers assess texture and shape; dorsal surfaces assess temperature; ulnar borders assess tactile fremitus.

\medterm{Pancreatic Adenocarcinoma}-- A highly lethal malignancy arising from the pancreatic ductal epithelium, with a five-year survival rate of less than 10 percent, making it one of the deadliest cancers. The most common risk factors include smoking, obesity, chronic pancreatitis, diabetes mellitus, and hereditary syndromes including BRCA2 mutations and Lynch syndrome. The pathogenesis involves sequential genetic mutations including KRAS activation, TP53 loss, SMAD4 loss, and CDKN2A inactivation. The majority of tumors arise in the head of the pancreas, presenting with painless obstructive jaundice. CA 19-9 is a useful serum tumor marker, though not specific for screening. Surgical resection (pancreaticoduodenectomy, Whipple procedure) offers the only chance of cure, but less than 20 percent of patients are resectable at presentation.

\medterm{Pancreatic Pseudocyst}-- A pancreatic pseudocyst is a localized collection of pancreatic fluid surrounded by a wall of granulation tissue without an epithelial lining, typically occurring as a complication of acute or chronic pancreatitis. Unlike true cysts, pseudocysts lack an epithelial lining, which is why they are called pseudocysts. The pathogenesis involves disruption of the pancreatic duct with leakage of pancreatic enzymes, causing enzymatic digestion of surrounding tissues and formation of a fluid collection that may reach large sizes. Pseudocysts present with epigastric pain, fullness, and tenderness. Complications include infection, rupture, and hemorrhage. Diagnosis is confirmed by CT or MRI showing a thick-walled fluid collection. Treatment options include observation, endoscopic drainage, or surgical cystogastrostomy.

\medterm{Parkinson's Disease}-- A progressive neurodegenerative disorder caused by selective loss of dopaminergic neurons in the substantia nigra pars compacta, leading to dopamine depletion in the striatum. The pathological hallmark is the presence of Lewy bodies, intracellular inclusions composed of aggregated alpha-synuclein protein, within surviving neurons. The disease typically manifests after age 60 with a male predominance. Clinical features include the cardinal motor symptoms of resting tremor, typically a pill-rolling tremor, bradykinesia, rigidity, and postural instability, as well as non-motor symptoms including depression, dementia, and autonomic dysfunction. Treatment is based on dopamine replacement with levodopa-carbidopa, dopamine agonists, and deep brain stimulation for advanced disease with motor fluctuations.

\medterm{Pathologic Hypertrophy}-- The enlargement of an organ resulting from an increase in the size of its constituent cells without an increase in cell number, typically occurring as a response to chronically increased workload or hormonal stimulation. Unlike physiologic hypertrophy, pathologic hypertrophy occurs under abnormal conditions and may eventually decompensate. The most important example is cardiac hypertrophy in response to chronic pressure overload from systemic hypertension or aortic stenosis. Initially compensatory and beneficial, sustained pathologic hypertrophy leads to maladaptive changes including myocardial fibrosis, diastolic dysfunction, and eventual systolic failure. The transition from compensated hypertrophy to heart failure involves altered calcium handling, energy metabolism, and apoptosis of cardiomyocytes.

\medterm{Petechiae}-- Small, pinpoint, non-blanching red or purple spots on the skin or mucous membranes caused by extravasation of red blood cells from capillaries into the surrounding tissue. They are typically less than 2 millimeters in diameter and result from decreased platelet count, platelet dysfunction, or increased vascular fragility. The most common cause is thrombocytopenia, where platelet counts fall below 20,000 to 30,000 per microliter, impeding primary hemostasis. Other causes include vascular fragility from vasculitis, connective tissue disorders, or corticosteroid use. Petechiae are non-blanching and extremely small. Their distribution provides diagnostic clues: dependent areas suggest thrombocytopenia, while distribution above the nipple line suggests superior vena cava syndrome. Diagnosis involves evaluating platelet count and coagulation studies.

\medterm{Platelet Plug Formation}-- The initial hemostatic response to vascular injury, involving platelet adhesion, activation, and aggregation. When the endothelium is disrupted, subendothelial von Willebrand factor and collagen are exposed to flowing blood. Platelet glycoprotein Ib-IX-V complex binds to vWF under high shear stress, tethering and decelerating platelets on the damaged surface. Collagen binding to glycoprotein VI and integrin alpha-2-beta-1 triggers intracellular signaling cascades through activation of phospholipase C and increased cytosolic calcium, leading to conformational change of integrin alpha-IIb-beta-3 which binds fibrinogen, cross-linking platelets into an aggregated plug. The plug is stabilised by the coagulation cascade generating a fibrin mesh. This primary hemostatic response is critical for stopping bleeding but can contribute to pathological thrombosis.

\medterm{Pressure Atrophy}-- The gradual wasting of tissue resulting from sustained external compression that impairs blood supply and causes cellular shrinkage and eventual death. The mechanism involves compression of blood vessels, which reduces perfusion and causes ischemia of the dependent tissue. Prolonged ischemia leads to cellular atrophy with reduced cell size and organ function, and if compression continues, progresses to necrosis. Clinical examples include decubitus ulcers that develop over bony prominences in bedridden patients, and atrophy of normal tissue adjacent to an expanding tumor. Affected cells show reduced size, decreased protein synthesis, and increased autophagic vacuoles. Unlike necrosis, pressure atrophy is a reversible adaptive change if the compressive force is removed before irreversible cell damage occurs.

\medterm{Prion Disease}-- Prion diseases, also known as transmissible spongiform encephalopathies, are a group of invariably fatal neurodegenerative disorders caused by misfolding of the normal cellular prion protein PrPC into the abnormal scrapie form PrPSc. The misfolded protein acts as an infectious agent by inducing normal PrPC to adopt the pathological conformation, propagating the misfolding process through an autocatalytic mechanism. PrPSc accumulates in the brain, causing neuronal death and characteristic spongiform change in the brain. The most common human prion disease is sporadic Creutzfeldt-Jakob disease, presenting with rapidly progressive dementia, myoclonus, and ataxia. Variant CJD is linked to bovine spongiform encephalopathy. Diagnosis is based on clinical features, MRI, EEG, and detection of 14-3-3 protein in CSF. All prion diseases are uniformly fatal.

\medterm{Prostate Adenocarcinoma}-- The most common non-cutaneous malignancy in men and the second leading cause of cancer death in men after lung cancer. The disease typically arises in the peripheral zone of the prostate and is hormone-dependent, requiring androgens for growth. Risk factors include age, family history, and race, with African American men having higher incidence and mortality. Most tumors are slow-growing and may never cause symptoms, while a subset are aggressive and metastasize. The Gleason grading system, based on histological glandular architecture, is the most important prognostic factor, with scores ranging from 6 to 10. PSA screening has improved early detection but remains controversial due to overdiagnosis of indolent disease. Treatment options include active surveillance, radical prostatectomy, radiation therapy, and androgen deprivation therapy.

\medterm{Purpura}-- Purple discoloration of the skin or mucous membranes caused by extravasation of red blood cells from blood vessels into the surrounding tissue. Purpura can be classified by size: petechiae are less than 2 millimeters, purpura are 2 to 10 millimeters, and ecchymoses are greater than 10 millimeters. Causes include thrombocytopenia from decreased production or increased destruction, platelet dysfunction from medications or uremia, vascular fragility from aging or vitamin C deficiency (scurvy), and Henoch-Schonlein purpura with IgA immune complex deposition. Palpable purpura suggests underlying vasculitis. Diagnosis involves complete blood count, coagulation studies, and skin biopsy if vasculitis is suspected. Treatment targets the underlying cause; corticosteroids are used in vasculitic causes.

\medterm{Pyelonephritis}-- An infection of the renal parenchyma and pelvicalyceal system, typically resulting from ascending urinary tract infection. Acute pyelonephritis most commonly affects women and is caused by gram-negative enteric organisms, particularly Escherichia coli, which account for approximately 80 percent of cases. The pathogenesis involves bacterial ascent from the bladder to the kidneys, facilitated by vesicoureteral reflux, urinary obstruction, and bacterial virulence factors including adhesins and toxins. Acute pyelonephritis presents with fever, chills, flank pain, nausea, and vomiting, with pyuria and bacteriuria on urinalysis. Treatment involves antibiotics guided by culture sensitivity. Complications include renal abscess, emphysematous pyelonephritis, sepsis, and chronic pyelonephritis leading to renal scarring and hypertension.

\medterm{Red Infarction}-- Red infarction, also known as hemorrhagic infarction, occurs in tissues with dual blood supply, loose tissue architecture, or venous obstruction, allowing blood to accumulate in the necrotic area and produce a dark red appearance. The hemorrhagic appearance results from extravasation of red blood cells from congested vessels into the necrotic tissue. Red infarcts characteristically occur in the lung, where the dual pulmonary arterial and bronchial arterial circulation provides collateral flow, allowing blood to seep into the necrotic area from adjacent patent vessels. In the liver and small intestine, red infarction can also occur due to dual blood supply. The affected tissue is dark red, wedge-shaped, and poorly demarcated. Organisation and scar formation follow, with functional impact depending on the size and location of the infarct.

\medterm{Renal Infarction}-- Renal infarction results from occlusion of the renal artery or its branches, causing coagulative necrosis of the renal parenchyma. The kidney has end-arterial supply with limited collateral circulation, making it susceptible to ischemic injury even from small vessel occlusion. Common causes include thromboembolism from the heart, aortic dissection extending into the renal arteries, atherosclerotic disease, and fibromuscular dysplasia. Clinical features include sudden onset severe flank pain, haematuria, nausea, and vomiting. Hypertension may develop due to renin release from the ischemic kidney. Diagnosis is confirmed by CT angiography or renal scintigraphy. Treatment depends on the cause and may include anticoagulation, revascularisation, or conservative management with analgesics and blood pressure control.

\medterm{Sarcoma}-- See oncology.

\medterm{Septic Embolism}-- Septic embolism occurs when infected material, typically thrombi containing bacteria, embolizes through the bloodstream to distant sites, causing metastatic infection. The most common source is infective endocarditis, where infected vegetations on heart valves break off and lodge in peripheral vessels. Other sources include infected central venous catheters, septic pelvic thrombophlebitis, and infected arterial thrombi. Septic emboli typically lodge in the lungs, causing multiple bilateral wedge-shaped pulmonary infarcts with cavitation, presenting with chest pain, haemoptysis, and fever. Systemic embolisation can cause abscesses in the brain, spleen, kidneys, and skin. Diagnosis involves blood cultures and imaging. Treatment requires prolonged antibiotics and often surgical intervention for source control.

\medterm{Septic Shock}-- A form of distributive shock caused by overwhelming infection and the host dysregulated immune response, leading to systemic vasodilation, increased vascular permeability, myocardial depression, and multiorgan dysfunction. It is the most common cause of death in intensive care units worldwide. The pathogenesis involves recognition of microbial products including lipopolysaccharide and peptidoglycan by pattern recognition receptors on immune cells, triggering massive release of pro-inflammatory cytokines including TNF-alpha, IL-1, and IL-6, causing widespread vasodilation, increased capillary permeability, and myocardial depression. This results in refractory hypotension, tissue hypoperfusion, and multiorgan failure. Management includes early broad-spectrum antibiotics, fluid resuscitation, vasopressors, and source control of infection.

\medterm{Shock}-- See emergency_medicine.

\medterm{Splenic Infarction}-- Splenic infarction occurs when the splenic artery or its branches are occluded, typically by emboli from the heart. The spleen has end-arterial supply with limited collateral circulation, making it vulnerable to infarction even with small vessel occlusion. The most common cause is cardioembolism, particularly in patients with atrial fibrillation, infective endocarditis, or left ventricular thrombus. Other causes include hematologic disorders like sickle cell disease and myeloproliferative disorders such as polycythaemia vera. Splenic infarction presents with left upper quadrant pain, often pleuritic, sometimes radiating to the left shoulder. Fever and splenomegaly may be present. Most cases are managed conservatively with analgesics and treatment of the underlying cause, though complications include abscess formation and splenic rupture.

\medterm{Starling Forces}-- Starling forces describe the balance of hydrostatic and oncotic pressures that govern fluid movement across capillary walls, determining the distribution of fluid between the intravascular and extravascular compartments. The Starling equation states that net filtration equals the filtration coefficient times the difference between capillary and interstitial hydrostatic pressures minus the difference between capillary and interstitial oncotic pressures. Capillary hydrostatic pressure, generated by arterial blood pressure, pushes fluid out of the capillary at the arteriolar end. Plasma oncotic pressure, generated primarily by albumin, pulls fluid back into the capillary at the venular end. The net result is that approximately 90 percent of filtered fluid is reabsorbed; the remainder drains via the lymphatic system as lymph.

\medterm{Storage Diseases}-- A group of inherited metabolic disorders characterized by accumulation of substrates within cells due to specific enzyme deficiencies. Lysosomal storage diseases result from defects in lysosomal hydrolytic enzymes, leading to accumulation of undigested macromolecules within lysosomes. Gaucher disease results from glucocerebrosidase deficiency, causing glucocerebroside accumulation in macrophages that become enlarged Gaucher cells with characteristic wrinkled tissue paper cytoplasm. Niemann-Pick disease results from sphingomyelinase deficiency causing sphingomyelin accumulation in macrophages. Tay-Sachs disease results from hexosaminidase A deficiency causing GM2 ganglioside accumulation in neurons, leading to neurodegeneration. Mucopolysaccharidoses result from deficiencies of enzymes degrading glycosaminoglycans.

\medterm{Synaptic Transmission}-- The fundamental process by which neurons communicate with each other through specialized junctions called synapses. At chemical synapses, an action potential arriving at the presynaptic terminal opens voltage-gated calcium channels, allowing calcium influx that triggers fusion of synaptic vesicles with the presynaptic membrane through SNARE protein complexes. Neurotransmitters are released into the synaptic cleft by exocytosis and bind to specific receptors on the postsynaptic membrane, generating either excitatory postsynaptic potentials (EPSPs) or inhibitory postsynaptic potentials (IPSPs). Signal termination occurs via neurotransmitter reuptake by specific transporters, enzymatic degradation in the synaptic cleft, or diffusion away from the synapse. This process forms the basis of all neural communication and information processing.

\medterm{Systemic Sclerosis}-- Systemic sclerosis, also known as scleroderma, is a chronic autoimmune disease characterized by progressive fibrosis of the skin and internal organs, vasculopathy, and immune dysfunction. The pathogenesis involves an aberrant wound healing response with excessive collagen production by activated fibroblasts, triggered by endothelial injury and immune activation. The disease exists in two main forms: limited cutaneous systemic sclerosis, formerly called CREST syndrome, with skin involvement distal to the elbows and knees, and diffuse cutaneous systemic sclerosis with widespread skin thickening and early internal organ involvement including interstitial lung disease, pulmonary hypertension, and renal crisis. Anti-centromere antibodies are associated with limited disease; anti-Scl-70 (anti-topoisomerase I) is associated with diffuse disease.

\medterm{Takayasu Arteritis}-- A chronic granulomatous vasculitis primarily affecting the aorta and its major branches, most commonly diagnosed in women under 50 years of age, particularly in Asian populations. The disease involves transmural inflammation of the aortic wall and its branches, leading to intimal fibrosis, luminal narrowing, and sometimes aneurysmal dilation. The pathogenesis involves inappropriate immune-mediated vascular inflammation with CD4 T cell and macrophage infiltration. Clinical features include absent or diminished pulses, blood pressure discrepancies between arms, vascular bruits, claudication, and syncope. Diagnosis is based on imaging showing vascular wall thickening, stenosis, or aneurysm. Treatment includes high-dose corticosteroids and other immunosuppressive agents to control inflammation and prevent disease progression.

\medterm{Temporal Arteritis}-- A vasculitic disorder affecting medium and large arteries, particularly the temporal artery, predominantly occurring in adults over 50 years of age. The disease is characterized by granulomatous inflammation of the arterial wall with fragmentation of the internal elastic lamina, intimal proliferation, and luminal narrowing that can lead to ischemia of affected tissues. The most feared complication is anterior ischemic optic neuropathy, which can cause sudden, irreversible vision loss. Other symptoms include headache, jaw claudication, scalp tenderness, and constitutional symptoms such as fever and weight loss. ESR and CRP are typically markedly elevated. Temporal artery biopsy shows granulomatous inflammation with giant cells and elastic lamina fragmentation. Prompt high-dose corticosteroid therapy protects against vision loss.

\medterm{Teratoma}-- See obstetrics_gynecology.

\medterm{Thrombosis}-- The formation of a blood clot within the cardiovascular system during life, potentially obstructing blood flow and causing tissue ischemia or infarction. Arterial thrombi typically form on ruptured atherosclerotic plaques, are composed primarily of platelets and fibrin, and appear pale and grossly attached to the vessel wall. Venous thrombi form in areas of stasis, are rich in fibrin and trapped red blood cells, and appear red with a characteristic Lines of Zahn layering pattern. Venous thrombi commonly cause deep vein thrombosis with leg swelling and pain, and may embolise to the lungs causing pulmonary embolism. Arterial thrombosis leads to myocardial infarction, stroke, or peripheral arterial occlusion. Risk factors are described by Virchow's triad. Treatment involves anticoagulants, thrombolytics, or mechanical thrombectomy.

\medterm{Thrombus Organization}-- The process by which a thrombus is gradually replaced by granulation tissue and eventually fibrous scar tissue. The process begins when inflammatory cells, particularly macrophages, infiltrate the thrombus and begin digesting the fibrin and dead cells through fibrinolysis. Fibroblasts then migrate into the thrombus and begin producing collagen and other extracellular matrix components, gradually replacing the original clot. New blood vessels grow into the thrombus through the process of angiogenesis, a phenomenon called recanalisation. The organized thrombus becomes a fibrous scar that may contract and narrow the vessel lumen or, in the case of recanalisation, restore partial blood flow. This process is part of the natural healing response to thrombosis.

\medterm{Thyroid Follicular Carcinoma}-- approximately 10 to 15 percent of thyroid cancers and is the second most common differentiated thyroid malignancy. Unlike papillary carcinoma, follicular carcinoma typically spreads hematogenously to distant sites including bone, lung, and liver, rather than to regional lymph nodes. The diagnosis requires demonstration of capsular invasion or vascular invasion on histological examination, distinguishing it from follicular adenoma which lacks these features. Prognosis is generally good, with ten-year survival exceeding 85 percent, though worse in the presence of distant metastasis. Treatment includes total thyroidectomy followed by radioactive iodine ablation and TSH suppression therapy.

\medterm{Thyroid Papillary Carcinoma}-- The most common type of thyroid cancer, accounting for approximately 80 percent of all thyroid malignancies, and has an excellent prognosis with over 95 percent 10-year survival. The tumor arises from follicular thyroid cells and is characterized by papillary architectural growth pattern and characteristic nuclear features including nuclear grooves, intranuclear inclusions, and orphan eye nuclei. Risk factors include childhood radiation exposure, particularly from therapeutic radiation to the head and neck in childhood. The tumor has a propensity for lymphatic spread, with cervical lymph node metastases present in up to 50 percent of cases at diagnosis. However, the prognosis remains excellent. Treatment involves thyroidectomy with or without radioactive iodine ablation and TSH suppression.

\medterm{TNF Alpha}-- Tumor necrosis factor alpha is a potent pro-inflammatory cytokine produced primarily by activated macrophages, as well as by T cells, NK cells, mast cells, and endothelial cells. TNF-alpha exists as a 26-kilodalton transmembrane protein that is cleaved by the metalloprotease TACE to release a soluble 17-kilodalton active form. It acts through two receptors: TNFR1, which contains a death domain and can activate both NF-kappaB survival signaling and caspase-8-mediated apoptosis, and TNFR2, which lacks a death domain and preferentially activates NF-kappaB, promoting cell survival and inflammation. TNF-alpha is a key therapeutic target in chronic inflammatory diseases including rheumatoid arthritis, psoriasis, and inflammatory bowel disease, where monoclonal antibodies (infliximab, adalimumab) and fusion proteins (etanercept) are used.

\medterm{Transudate}-- A transudate is a type of edema fluid with low protein content and few cells, resulting from imbalances in hydrostatic and oncotic pressures across capillary walls rather than from inflammation or increased vascular permeability. Transudates are typically caused by conditions that increase capillary hydrostatic pressure, such as congestive heart failure causing pulmonary edema, or decrease plasma oncotic pressure, such as hypoalbuminemia from liver cirrhosis or nephrotic syndrome. The fluid is clear, straw-coloured, and has low cellularity. Laboratory criteria for transudate include pleural fluid protein-to-serum protein ratio less than 0.5, pleural fluid LDH-to-serum LDH ratio less than 0.6, and pleural fluid LDH less than two-thirds the upper limit of normal for serum.

\medterm{Tumor Differentiation}-- The degree to which neoplastic cells morphologically and functionally resemble their normal counterparts of origin. Well-differentiated tumors retain many structural features of the original tissue, including organized architectural patterns, specialized cellular structures, and sometimes functional capabilities such as hormone production or mucin secretion. Poorly differentiated or anaplastic tumors show minimal resemblance to the normal tissue and often require immunohistochemistry and molecular profiling for accurate classification. tumor grade, a measure of differentiation, is a powerful independent prognostic factor in most cancers. Well-differentiated tumors generally have a better prognosis than poorly differentiated ones, though notable exceptions exist.

\medterm{Tumor Embolism}-- Tumor embolism occurs when malignant cells, tumor fragments, or tumor thrombi enter the venous circulation and travel to distant sites. The pulmonary vasculature is the most common site of embolization, as venous blood from most organs drains to the lungs through the pulmonary arteries. Tumor emboli can also reach the arterial circulation through pulmonary capillary shunts or a patent foramen ovale, causing systemic embolization to the brain, kidneys, or other organs. The clinical presentation ranges from asymptomatic pulmonary emboli to acute respiratory failure and cor pulmonale. tumor embolism is an under-recognised complication of malignancy and may mimic thromboembolic disease. Diagnosis requires a high index of suspicion and histopathological confirmation.

\medterm{Tumor Suppressor Genes}-- Tumor suppressor genes encode proteins that normally inhibit cell proliferation, promote DNA repair, or induce programmed cell death in response to cellular damage. They function as brakes on the cell cycle, and loss of their function removes critical checkpoints that normally prevent uncontrolled growth. According to Knudson's two-hit hypothesis, both alleles must be inactivated for loss of function, which can occur through mutation, deletion, or epigenetic silencing. Hereditary cancer syndromes, such as Li-Fraumeni syndrome (TP53), familial adenomatous polyposis (APC), and hereditary breast and ovarian cancer (BRCA1, BRCA2), involve germline mutations of tumor suppressor genes. Epigenetic silencing through promoter hypermethylation is another common mechanism of inactivation in sporadic cancers.

\medterm{Urothelial Carcinoma}-- Urothelial carcinoma, previously called transitional cell carcinoma, is the most common malignancy of the urinary tract, arising from the urothelium that lines the renal pelvis, ureters, bladder, and proximal urethra. The most important risk factor is cigarette smoking, which accounts for approximately 50 percent of cases, with other risk factors including occupational exposure to aromatic amines, cyclophosphamide therapy, and chronic schistosomiasis. The pathogenesis involves field cancerization, where the entire urothelium is exposed to carcinogens, leading to multifocal and recurrent tumors. Gross haematuria is the most common presenting symptom. Diagnosis is by cystoscopy and biopsy. Treatment depends on stage: transurethral resection and intravesical therapy for non-muscle-invasive disease, and radical cystectomy for muscle-invasive disease.

\medterm{Valvular Heart Disease}-- Valvular heart disease involves dysfunction of one or more heart valves, leading to stenosis, regurgitation, or a combination of both. The most common causes are degenerative calcification of aortic and mitral valves, rheumatic heart disease, infective endocarditis, and ischemic papillary muscle dysfunction. Aortic stenosis results from calcification of a congenitally bicuspid or tricuspid aortic valve, causing left ventricular outflow obstruction with pressure overload. Aortic regurgitation results from incomplete valve closure, causing volume overload and left ventricular dilation. Mitral stenosis is most commonly rheumatic in origin, causing left atrial enlargement and pulmonary hypertension. Treatment includes medical management and valve replacement or repair for severe symptomatic disease.

\medterm{Vasodilation}-- The widening of blood vessels during inflammation, resulting in increased blood flow to the affected area. This process is mediated primarily by histamine released from mast cell granules acting on H1 receptors, nitric oxide produced by endothelial nitric oxide synthase, and prostaglandins particularly PGE2. The increased blood flow is responsible for the redness and heat characteristic of acute inflammation. Vasodilation occurs within seconds to minutes of tissue injury and represents the earliest vascular response in the inflammatory cascade. The increased blood flow brings oxygen, nutrients, and immune cells to the affected tissue and dilutes toxins. Vasodilation is accompanied by increased vascular permeability, allowing exudation of protein-rich fluid and formation of tissue edema, contributing to the cardinal signs of inflammation.

\medterm{Vegetation}-- A vegetation is a mass composed of platelets, fibrin, and microorganisms that forms on a heart valve or endocardial surface in the setting of infective endocarditis. The vegetation typically attaches to the valve leaflet on the low-pressure side and may progressively destroy the underlying valve tissue, causing valvular regurgitation. The most common causative organisms include Staphylococcus aureus in acute endocarditis, which can affect normal valves, and viridans streptococci in subacute endocarditis, which typically affects previously damaged valves. Diagnosis relies on blood cultures and echocardiography, with transoesophageal echocardiography being more sensitive. Treatment requires prolonged intravenous antibiotics and sometimes surgical valve debridement or replacement.

\medterm{Virchow's Triad}-- Virchow's triad describes three broad categories of factors that contribute to thrombosis: endothelial injury or dysfunction, stasis or turbulent blood flow, and hypercoagulability. Endothelial injury exposes subendothelial collagen and tissue factor, initiating platelet adhesion and the coagulation cascade. Causes include atherosclerosis, hypertension, vasculitis, and direct trauma. Stasis or turbulent flow disrupts laminar blood flow, bringing platelets into contact with the endothelium, preventing dilution of activated clotting factors and causing endothelial hypoxia. Hypercoagulability encompasses inherited conditions such as factor V Leiden, prothrombin gene mutation, and deficiencies of protein C, protein S, or antithrombin III, as well as acquired states including pregnancy, malignancy, and oral contraceptive use.

\medterm{Wegener's Granulomatosis}-- Wegener's granulomatosis, now called granulomatosis with polyangiitis, is a systemic vasculitis characterized by necrotizing granulomatous inflammation of the upper and lower respiratory tracts and pauci-immune necrotizing glomerulonephritis. It is associated with anti-neutrophil cytoplasmic antibodies, specifically c-ANCA targeting proteinase 3. The disease affects small to medium-sized vessels and can involve virtually any organ system. Upper respiratory tract involvement causes chronic sinusitis, nasal mucosal ulceration, and saddle nose deformity. Pulmonary involvement causes cough, haemoptysis, and cavitating lung nodules. Renal involvement, if untreated, rapidly progresses to end-stage renal disease. Treatment includes high-dose corticosteroids with cyclophosphamide or rituximab, which has dramatically improved outcomes.

\medterm{White Infarction}-- White infarction, also known as anemic infarction, occurs in solid organs with end-arterial circulation where collateral blood supply is limited, resulting in pale, well-demarcated areas of coagulative necrosis. The affected tissue appears pale because compression of the capillary bed by interstitial edema limits secondary hemorrhage into the necrotic area. White infarcts characteristically occur in the heart, kidney, spleen, and other solid organs with single-source arterial supply. The mechanism involves sudden arterial occlusion producing a wedge-shaped area of coagulative necrosis, with the apex pointing toward the occlusion site. The infarct appears pale because the necrotic tissue loses its blood content and surrounding vessels constrict. Organisation and fibrous scar formation follow over several weeks.

\medterm{Xeroderma Pigmentosum}-- A rare autosomal recessive disorder of nucleotide excision repair (NER) characterized by extreme sun sensitivity and a 1000-fold increased risk of skin cancers (basal cell carcinoma, squamous cell carcinoma, and melanoma, often in childhood). Mutations occur in one of eight genes (XPA-XPG and XP-variant/polymerase eta). Defective NER leads to accumulation of UV-induced pyrimidine dimers and other photoproducts, causing mutations in tumor suppressor genes (p53, PTCH). Clinical features: severe sunburn with minimal sun exposure, poikiloderma (hyperpigmentation, hypopigmentation, telangiectasias, atrophy), xerosis (dry skin), actinic keratoses, and multiple skin cancers. Ocular involvement (photophobia, conjunctivitis, keratitis, eyelid skin cancers) and neurological involvement (progressive neurodegeneration in XP-A, -B, -D, -G: microcephaly, spasticity, chorea, sensorineural hearing loss, peripheral neuropathy, cognitive impairment). Diagnosis: clinical criteria with confirmatory DNA repair assay (unscheduled DNA synthesis/UDS) showing reduced NER capacity or gene sequencing. Management: rigorous sun protection (sunscreen SPF 50+, UV-protective clothing, hats, sunglasses), regular dermatological surveillance (every 3-6 months) with early treatment of skin cancers (cryotherapy, photodynamic therapy, surgical excision, topical 5-fluorouracil, imiquimod). Oral retinoids (acitretin) reduce skin cancer incidence by 50-70\% in XP patients. Prognosis: reduced life expectancy primarily from metastatic skin cancer or progressive neurodegeneration. Neurological complications are the major determinant of poor prognosis.
